% !TEX program = pdflatex
\documentclass[11pt]{article}
\usepackage[margin=1in]{geometry}
\usepackage{pgfplots}
\pgfplotsset{compat=1.18}
\usepackage{amsmath}

\title{Example 1: Two-Body Orbit with Precession}
\author{WeberElectrodynamics.jl}
\date{\today}

\begin{document}
\maketitle

\section{Weber Trajectory}

Figure~\ref{fig:trajectory} shows a rosette pattern from Weber electrodynamics ($c = 4$).

\begin{figure}[htbp]
    \centering
    \input{figures/trajectory_weber.tikz}
    \caption{Two-body trajectory under Weber electrodynamics. The velocity-dependent potential causes apsidal precession, creating a rosette pattern.}
    \label{fig:trajectory}
\end{figure}

\section{Comparison with Coulomb Limit}

Figure~\ref{fig:comparison} compares Weber ($c = 4$) with the Coulomb limit ($c \to \infty$).

\begin{figure}[htbp]
    \centering
    \input{figures/trajectory_comparison.tikz}
    \caption{Weber (solid blue) vs Coulomb (dashed red). The Coulomb trajectory forms a closed ellipse; the Weber trajectory precesses.}
    \label{fig:comparison}
\end{figure}

\end{document}
